\subsection*{3.6. Aplicación a la descripción de la categoría cofibrada
de las biextensiones de una pareja dada $(P,Q)$ de grupos conmutativos}

Sean $P$ y $Q$ dos $\mathbb Z$-módulos del topos $\underline T$.
Elijamos una resolución plana $L'_\bullet(P)$ de $P$ que coincida en
grados a lo sumo $2$ con el complejo aumentado $L_\bullet(P)$ de 3.5.
Esto es posible por 3.5.3 y porque todo $\mathbb Z$-módulo sobre
$\underline T$ es cociente de un $\mathbb Z$-módulo libre y, por tanto,
plano. Elijamos análogamente una resolución plana $L'_\bullet(Q)$ de $Q$
que coincida con $L_\bullet(Q)$ en grados a lo sumo $2$. Se tienen
homomorfismos canónicos
\begin{equation}\tag{3.6.1}
 L_\bullet(P)\longrightarrow L'_\bullet(P),
 \qquad
 L_\bullet(Q)\longrightarrow L'_\bullet(Q),
\end{equation}
y, por consiguiente, un homomorfismo canónico
\begin{equation}\tag{3.6.2}
\begin{split}
 L_\bullet(P,Q)
 &\overset{\mathrm{def}}{=}L_\bullet(P)\otimes L_\bullet(Q)
 \longrightarrow L'_\bullet(P,Q)\\
 &\overset{\mathrm{def}}{=}L'_\bullet(P)\otimes L'_\bullet(Q)
 \simeq P\mathbin{\overset{\mathbb L}{\otimes}}Q.
\end{split}
\end{equation}
Es un isomorfismo sobre las componentes en grados a lo sumo $2$. Por
consiguiente, por definición,
\begin{equation}\tag{3.6.3}
 \underline\Psi_{L_\bullet(P,Q)}
 =
 \sisi{\underline\Psi_{L_\bullet(P',Q')}}%
       {\underline\Psi_{L'_\bullet(P,Q)}}.
\end{equation}

\underline{Teorema} 3.6.4. \underline{La categoría cofibrada aditiva
sobre la categoría}
\[
 \underline C=\operatorname{GRAB}(\underline T)
\]
\underline{de los grupos conmutativos de $\underline T$, formada por las
biextensiones de $(P,Q)$ por un grupo abeliano variable $G$} (2.4),
\underline{es canónicamente equivalente a la categoría cofibrada aditiva}
\[
 \underline\Psi_{L_\bullet(P,Q)}
 =
 \underline\Psi_{L'_\bullet(P,Q)}
\]
\underline{definida a partir de}
$L_\bullet(P,Q)=L_\bullet(P)\otimes L_\bullet(Q)$
\underline{como en 3.1}.

Esto y 3.2.5 proporcionan el resultado principal de la presente sección,
anunciado anteriormente:

%%%%%%%%%%  scan idx 216  |  folio 205  %%%%%%%%%%

\underline{Corolario} 3.6.5. \underline{Para $P$ y $Q$ fijos hay un
isomorfismo canónico, funtorial en $G$,}
\begin{equation}\tag{3.6.5.1}
 \Biext^1(P,Q;G)
 \xrightarrow{\sim}
 \Ext^1(P\mathbin{\overset{\mathbb L}{\otimes}}Q,G).
\end{equation}

Para demostrar 3.6.5, explicitemos primero las componentes de grados
$0,1,2$ de $L_\bullet(P,Q)$:
\begin{equation}\tag{3.6.6}
\left\{
\begin{aligned}
 L_0(P,Q)
  ={}&L_0(P)\otimes L_0(Q)
     =\mathbb Z[P\times Q],\\
 L_1(P,Q)
  ={}&L_0(P)\otimes L_1(Q)+L_1(P)\otimes L_0(Q)\\
  ={}&\mathbb Z[P\times Q\times Q]
     +\mathbb Z[P\times P\times Q],\\
 L_2(P,Q)
  ={}&L_0(P)\otimes L_2(Q)+L_1(P)\otimes L_1(Q)
     +L_2(P)\otimes L_0(Q)\\
  ={}&\mathbb Z[P\times Q\times Q]
     +\mathbb Z[P\times Q\times Q\times Q]\\
   &+\mathbb Z[P\times P\times Q\times Q]\\
   &+\mathbb Z[P\times P\times Q]
     +\mathbb Z[P\times P\times P\times Q].
\end{aligned}
\right.
\end{equation}

Para describir la categoría cofibrada
$\underline\Psi_{L_\bullet(P,Q)}$ mediante sus fibras
$\underline\Psi_{L_\bullet(P,Q)}(G)$, usamos 1.4 para expresar en términos
de torsores las extensiones de $L_i(P,Q)$, para $i=0,1,2$, por $G$. Una
extensión de $\sisi{L_0(P\times Q)}{L_0(P,Q)}$ por $G$ es un torsor bajo
$G_{P\times Q}$. Una extensión de $L_1(P,Q)$ por $G$ es una pareja formada
por un torsor bajo $G_{P\times Q\times Q}$ y otro bajo
$G_{P\times P\times Q}$. De manera análoga, una extensión de $L_2(P,Q)$
por $G$ es un sistema de cinco torsores bajo los grupos obtenidos de $G$
por cambio de base a
\[
\begin{gathered}
 P\times Q\times Q,\quad
 P\times Q\times Q\times Q,\quad
 P\times P\times Q\times Q,\\
 P\times P\times Q,\quad
 P\times P\times P\times Q.
\end{gathered}
\]

Bajo este diccionario, una pareja $(X,\alpha)$ formada por una extensión
de $L_0(P,Q)$ por $G$ y una trivialización de $d_0^*X$ sobre $L_1(P,Q)$
se convierte en un torsor $E$ sobre $P\times Q$, dotado de dos
trivializaciones $(\varphi,\psi)$ de los torsores $(E_1,E_2)$ sobre
$P\times Q\times Q$ y $P\times P\times Q$, respectivamente, obtenidos
de $E$ por imagen inversa mediante

%%%%%%%%%%  scan idx 217  |  folio 206  %%%%%%%%%%

\[
 d_0:L_1(P,Q)\longrightarrow L_0(P,Q).
\]
Al explicitar $d_0$, como en la demostración de 3.5.3, se identifican
$\varphi$ y $\psi$ con los isomorfismos
\[
 \varphi:q_1^*E\,q_2^*E\xrightarrow{\sim}\pi_Q^*E,
 \qquad
 \psi:p_1^*E\,p_2^*E\xrightarrow{\sim}\pi_P^*E.
\]
Aquí $\pi_Q,\sisi{q_1q_2}{q_1,q_2}:P\times Q\times Q
\longrightarrow\sisi{Q}{P\times Q}$ están inducidos por la ley de
composición y por las proyecciones primera y segunda
$Q\times Q\longrightarrow Q$; las aplicaciones
\[
 \pi_P,p_1,p_2:P\times P\times Q\longrightarrow P\times Q
\]
se definen simétricamente. Así, $\varphi$ y $\psi$ son datos de las formas
(2.0.4) y (2.0.7) que aparecen en 2.1.

Finalmente, la condición de que $(X,\alpha)$ sea un objeto de
$\underline\Psi_{L_\bullet(P,Q)}(G)$, a saber, la compatibilidad de la
trivialización parcial $\alpha$ con $d_0d_1=0$, se convierte en cinco
relaciones de compatibilidad para la estructura $(E,\varphi,\psi)$ sobre
los torsores definidos, respectivamente, sobre
\[
\begin{gathered}
 P\times Q\times Q,\quad
 P\times Q\times Q\times Q,\quad
 P\times P\times Q\times Q,\\
 P\times P\times Q,\quad
 P\times P\times P\times Q.
\end{gathered}
\]
Estas corresponden a los cinco términos de la expresión (3.6.6) para
$L_2(P,Q)$. Como en la demostración de 3.5.3, las dos primeras son
precisamente la relación de conmutatividad (2.0.6) y la relación de
asociatividad (2.0.5); juntas dicen que $\varphi$ dota a $E$ de la
estructura de una extensión conmutativa de $Q_P$ por $G_P$.
Simétricamente, las dos últimas dicen que $\psi$ dota a $E$ de la
estructura de una extensión conmutativa de $P_Q$ por
$\sisi{G_P}{G_Q}$. La relación intermedia, sobre los torsores definidos
sobre $P\times P\times Q\times Q$, es exactamente la conmutatividad del
diagrama (2.1.1). Dice que $\varphi$ y $\psi$ son compatibles y, en vista
de las otras condiciones, definen sobre $E$ una estructura de biextensión
de $(P,Q)$ por $G$. La comprobación de estas interpretaciones es inmediata
y se deja al lector.

%%%%%%%%%%  scan idx 218  |  folio 207  %%%%%%%%%%

La discusión precedente define efectivamente una equivalencia de
categorías cofibradas sobre
$\underline C=\operatorname{GRAB}(\underline T)$, tal como se afirmó en
3.6.4.

\underline{Observación} 3.6.7. Cada una de las tres generalizaciones de
las biextensiones mencionadas en 2.10 tiene una variante correspondiente
de 3.6.4 y, por tanto, de 3.6.5. La generalización a un número cualquiera
de factores (2.10.2) es inmediata, al igual que la generalización a las
biextensiones de $A$-módulos (2.10.3), una vez construida la variante de
$L_\bullet(P)$ descrita en 3.5.4 b). Desde el punto de vista del álgebra
homológica es más interesante el caso en que no se supone que $P$ y $Q$
sean conmutativos (2.10.1).

Asociemos a un grupo no necesariamente conmutativo $P$ de $\underline T$
un complejo $K_\bullet(P)$ truncado en grado $2$. Se describe como se
describió $L_\bullet(P)$ en 3.5, salvo que
\[
 K_2(P)=\mathbb Z[P\times P\times P],
\]
omitiéndose el primer factor que aparecía en $L_2(P)$; $d_0$ y $d_1$ se
definen como en (3.5.2). Entonces
\[
 H_0(K_\bullet(P))=P/[P,P]
 \simeq\underline H_1(P,\mathbb Z),
\]
y, análogamente,
\[
 H_1(K_\bullet(P))\simeq\underline H_2(P,\mathbb Z).
\]
En efecto, $K_\bullet(P)$ es precisamente el truncamiento en grado $2$ del
complejo de cadenas $C_\bullet(P,\mathbb Z)$ de $P$ con coeficientes en
$\mathbb Z$.

El argumento que demuestra 3.5.3 muestra asimismo que la categoría
cofibrada $\underline\Psi_{K_\bullet(P)}$ sobre
$\underline C=\operatorname{GRAB}(\underline T)$ es canónicamente

%%%%%%%%%%  scan idx 219  |  folio 208  %%%%%%%%%%

equivalente a la categoría de las extensiones centrales de $P$ por grupos
conmutativos variables $G$. Análogamente, para dos grupos $P,Q$, la
categoría cofibrada de las biextensiones de $(P,Q)$ en el sentido de
2.10.1 es equivalente a $\underline\Psi_{K_\bullet(P,Q)}$, donde
\[
 K_\bullet(P,Q)
 =
 K_\bullet(P)\mathbin{\overset{\mathbb L}{\otimes}}K_\bullet(Q).
\]

Estas consideraciones sugieren que, cuando $P$ es conmutativo, la
resolución canónica buscada en 3.5.4 a) podría obtenerse mediante una
modificación adecuada de $C_\bullet(P,\mathbb Z)$.
